% Generated draft: requires mathematical and source review.
\documentclass[11pt]{article}
\usepackage[margin=25mm]{geometry}
\usepackage{fontspec}
\setmainfont{Latin Modern Roman}
\usepackage{amsmath,amssymb,mathtools,unicode-math}
\setmathfont{Latin Modern Math}
\usepackage{enumitem,needspace,xurl,hyperref}
\setlength{\parindent}{0pt}
\setlength{\parskip}{5pt}
\setlength{\emergencystretch}{3em}
\newcommand{\N}{\mathbb{N}}
\newcommand{\Z}{\mathbb{Z}}
\newcommand{\Q}{\mathbb{Q}}
\newcommand{\R}{\mathbb{R}}
\newcommand{\C}{\mathbb{C}}
\newcommand{\F}{\mathbb{F}}
\newcommand{\A}{\mathbb{A}}
\newcommand{\PP}{\mathbb{P}}
\newcommand{\E}{\mathbb{E}}
\newcommand{\eps}{\varepsilon}
\newcommand{\dd}{\,\mathrm{d}}
\begin{document}
\section*{UnsolvedMath ID: 800302}
\textbf{Source problem code:} catalogue entry
\subsection{Definitions and mathematical statement}
Write \(S^1=\R/\Z\) and \(D=\Z[1/2]/\Z\). Thompson's group \(T\) is the group of orientation-preserving piecewise-linear homeomorphisms of \(S^1\) with finitely many breakpoints, all in \(D\), slopes in \(\{2^k:k\in\Z\}\), and preserving \(D\). Use its natural action on \(S^1\).

Let \(\mu:T\to[0,1]\) be a probability distribution satisfying
\[
\sum_{g\in T}\mu(g)=1,\qquad
\mu(g)=\mu(g^{-1}),\qquad
|\operatorname{supp}\mu|<\infty.
\]
Interpret the source's term \emph{non-degenerate} as requiring \(\langle\operatorname{supp}\mu\rangle=T\); the supplied record does not explicitly define this term. A Borel probability measure \(\nu\) on \(S^1\) is \(\mu\)-stationary if
\[
\nu=\sum_{g\in T}\mu(g)\,g_*\nu,
\qquad (g_*\nu)(A)=\nu(g^{-1}A)
\]
for every Borel set \(A\subseteq S^1\).

\textbf{Question (Navas, 2017).} For every symmetric, finitely supported, non-degenerate \(\mu\) as above, does \(S^1\), endowed with the corresponding \(\mu\)-stationary measure \(\nu\), realize the Poisson--Furstenberg boundary of \((T,\mu)\) as a measured \(T\)-space?

More explicitly, with the convention that a bounded real-valued function \(h:T\to\R\) is \(\mu\)-harmonic when
\[
h(g)=\sum_{s\in T}\mu(s)h(gs)\qquad(g\in T),
\]
the requested boundary identification means that the Poisson transform
\[
\mathcal P_\nu:L^\infty(S^1,\nu;\R)\longrightarrow H^\infty(T,\mu),
\qquad
(\mathcal P_\nu f)(g)=\int_{S^1}f(gx)\,\dd\nu(x),
\]
is an isometric bijection, where \(H^\infty(T,\mu)\) denotes the bounded real-valued \(\mu\)-harmonic functions with the supremum norm. Thus each such harmonic function must have a unique representing function modulo \(\nu\)-null sets, with equal norms.

\textbf{Unresolved qualification.} The source refers to ``the corresponding stationary measure'' without specifying its selection or asserting uniqueness. The question above retains that qualification; the supplied record alone does not determine whether a particular canonical measure, some stationary measure, or every stationary measure is intended.
\subsection{Short English statement}
For each symmetric, finitely supported probability distribution whose support generates Thompson's group \(T\), Navas asks whether the natural circle action with its corresponding stationary measure gives the full Poisson--Furstenberg boundary. This would represent every bounded harmonic function uniquely by a bounded measurable function on the circle. The intended selection of stationary measure remains unspecified in the supplied record.
\subsection{Sources}
\begin{itemize}
\item Local UnsolvedMath registry: supplied complete source record, permanent ID \(800302\), source problem code \texttt{catalogue entry}, upstream revision \texttt{d371ee917ac9647fbb06cd6fc529deec6f3dda82}.
\item Navas, \emph{Group Actions on 1-Manifolds: A List of Very Concrete Open Questions} (2017), source item recorded as ``Question 19 (\#18)''. Supplied source URL: \url{https://arxiv.org/abs/1712.06462}. The record attributes the problem to Navas and gives the proposed year as 2017.
\item Supplied provenance reports access on \texttt{2026-10-16} and extraction by \texttt{arxiv\_latex}. These are supplied provenance claims, not access or source verification performed for this draft. The registry labels the problem \emph{open}; that label is supplied curation, not an independently verified literature assessment.
\end{itemize}
\end{document}
